
\documentclass{article}
\usepackage{colortbl}
\usepackage{makecell}
\usepackage{multirow}
\usepackage{supertabular}

\begin{document}

\newcounter{utterance}

\centering \large Interaction Transcript for game `dond', experiment `coop', episode 4 with nemotron{-}3{-}ultra.
\vspace{24pt}

{ \footnotesize  \setcounter{utterance}{1}
\setlength{\tabcolsep}{0pt}
\begin{supertabular}{c@{$\;$}|p{.15\linewidth}@{}p{.15\linewidth}p{.15\linewidth}p{.15\linewidth}p{.15\linewidth}p{.15\linewidth}}
    \# & $\;$A & \multicolumn{4}{c}{Game Master} & $\;\:$B\\
    \hline

    \theutterance \stepcounter{utterance}  
    & & \multicolumn{4}{p{0.6\linewidth}}{
        \cellcolor[rgb]{0.9,0.9,0.9}{
            \makecell[{{p{\linewidth}}}]{
                \texttt{\tiny{[P1$\langle$GM]}}
                \texttt{Vous jouez à un jeu de négociation coopérative dans lequel vous devez vous entendre avec un autre joueur sur la façon de diviser un ensemble d'objets entre vous.} \\
\\ 
\texttt{Règles:} \\
\texttt{(a) Vous et l'autre joueur recevez un ensemble d'objets. Chacun de vous reçoit également une fonction de valeur secrète, représentant la valeur que vous accordez à chaque type d'objet.} \\
\texttt{(b) Vous échangez des messages avec l'autre joueur pour vous entendre sur qui obtient quel objet. Chacun peut envoyer un maximum de 5 messages ou terminer le jeu tôt en faisant une proposition secrète à tout moment.} \\
\texttt{(c) On vous demande de soumettre une proposition secrète indiquant les objets que vous souhaitez, formatée entre crochets comme suit: "{[}Proposition: <nombre> <nom de l'objet>, <nombre> <nom de l'objet>, <...>{]}"} \\
\texttt{(d) Si vos propositions sont complémentaires, c'est{-}à{-}dire qu'il y a suffisamment d'objets pour satisfaire les deux propositions, chaque joueur reçoit un score basé sur la somme des valeurs des objets qu'il a reçus. Sinon, vous recevez tous les deux zéro point.} \\
\texttt{(e) Votre objectif est de maximiser la somme de votre score et du score de l'autre joueur.} \\
\\ 
\texttt{Commençons.} \\
\\ 
\texttt{L'ensemble des objets disponibles est:} \\
\\ 
\texttt{3 livres, 1 sac, 1 boulon} \\
\\ 
\texttt{Votre fonction de valeur secrète est:} \\
\\ 
\texttt{livre: 0, sac: 2, boulon: 8} \\
\\ 
\texttt{IMPORTANT: Vos messages, sauf s'il s'agit d'une proposition secrète, sont transmis directement à l'autre joueur, veuillez donc ne pas inclure de réponse aux règles ou de texte annonçant votre message. Pour faire une proposition secrète, utilisez le format indiqué. N'utilisez pas de crochets lorsque vous communiquez avec l'autre joueur, sinon cela sera interprété comme votre proposition secrète.} \\
            }
        }
    }
    & \\ \\

    \theutterance \stepcounter{utterance}  
    & \multicolumn{4}{p{0.6\linewidth}}{
        \cellcolor[rgb]{0.9,0.9,0.9}{
            \makecell[{{p{\linewidth}}}]{
                \texttt{\tiny{[P1$\rangle$GM]}}
                \texttt{Je vois qu'on a 3 livres, 1 sac et 1 boulon à partager. Ma priorité absolue est le boulon (il vaut 8 pour moi), le sac a une valeur modérée (2), et les livres ne m'apportent rien (0).} \\
\\ 
\texttt{Quelles sont vos priorités sur ces objets ?} \\
            }
        }
    }
    & & \\ \\

    \theutterance \stepcounter{utterance}  
    & & & \multicolumn{2}{p{0.3\linewidth}}{
        \cellcolor[rgb]{0.9,0.9,0.9}{
            \makecell[{{p{\linewidth}}}]{
                \texttt{\tiny{[GM$|$GM]}}
                \texttt{continue} \\
            }
        }
    }
    & & \\ \\

    \theutterance \stepcounter{utterance}  
    & & & \multicolumn{2}{p{0.3\linewidth}}{
        \cellcolor[rgb]{0.9,0.9,0.9}{
            \makecell[{{p{\linewidth}}}]{
                \texttt{\tiny{[GM$|$GM]}}
                \texttt{Je vois qu'on a 3 livres, 1 sac et 1 boulon à partager. Ma priorité absolue est le boulon (il vaut 8 pour moi), le sac a une valeur modérée (2), et les livres ne m'apportent rien (0).} \\
\\ 
\texttt{Quelles sont vos priorités sur ces objets ?} \\
            }
        }
    }
    & & \\ \\

    \theutterance \stepcounter{utterance}  
    & & \multicolumn{4}{p{0.6\linewidth}}{
        \cellcolor[rgb]{0.9,0.9,0.9}{
            \makecell[{{p{\linewidth}}}]{
                \texttt{\tiny{[GM$\rangle$P2]}}
                \texttt{Vous jouez à un jeu de négociation coopérative dans lequel vous devez vous entendre avec un autre joueur sur la façon de diviser un ensemble d'objets entre vous.} \\
\\ 
\texttt{Règles:} \\
\texttt{(a) Vous et l'autre joueur recevez un ensemble d'objets. Chacun de vous reçoit également une fonction de valeur secrète, représentant la valeur que vous accordez à chaque type d'objet.} \\
\texttt{(b) Vous échangez des messages avec l'autre joueur pour vous entendre sur qui obtient quel objet. Chacun peut envoyer un maximum de 5 messages ou terminer le jeu tôt en faisant une proposition secrète à tout moment.} \\
\texttt{(c) On vous demande de soumettre une proposition secrète indiquant les objets que vous souhaitez, formatée entre crochets comme suit: "{[}Proposition: <nombre> <nom de l'objet>, <nombre> <nom de l'objet>, <...>{]}"} \\
\texttt{(d) Si vos propositions sont complémentaires, c'est{-}à{-}dire qu'il y a suffisamment d'objets pour satisfaire les deux propositions, chaque joueur reçoit un score basé sur la somme des valeurs des objets qu'il a reçus. Sinon, vous recevez tous les deux zéro point.} \\
\texttt{(e) Votre objectif est de maximiser la somme de votre score et du score de l'autre joueur.} \\
\\ 
\texttt{Commençons.} \\
\\ 
\texttt{L'ensemble des objets disponibles est:} \\
\\ 
\texttt{3 livres, 1 sac, 1 boulon} \\
\\ 
\texttt{Votre fonction de valeur secrète est:} \\
\\ 
\texttt{livre: 1, sac: 2, boulon: 5} \\
\\ 
\texttt{IMPORTANT: Vos messages, sauf s'il s'agit d'une proposition secrète, sont transmis directement à l'autre joueur, veuillez donc ne pas inclure de réponse aux règles ou de texte annonçant votre message. Pour faire une proposition secrète, utilisez le format indiqué. N'utilisez pas de crochets lorsque vous communiquez avec l'autre joueur, sinon cela sera interprété comme votre proposition secrète.} \\
\\ 
\texttt{Je vois qu'on a 3 livres, 1 sac et 1 boulon à partager. Ma priorité absolue est le boulon (il vaut 8 pour moi), le sac a une valeur modérée (2), et les livres ne m'apportent rien (0).} \\
\\ 
\texttt{Quelles sont vos priorités sur ces objets ?} \\
            }
        }
    }
    & \\ \\

    \theutterance \stepcounter{utterance}  
    & & & \multicolumn{4}{p{0.6\linewidth}}{
        \cellcolor[rgb]{0.9,0.9,0.9}{
            \makecell[{{p{\linewidth}}}]{
                \texttt{\tiny{[GM$\langle$P2]}}
                \texttt{Je préfère largement le boulon (5 points pour moi), le sac a une valeur modérée (2 points), et les livres ont une valeur faible (1 point chacun). } \\
\\ 
\texttt{Vu que nous voulons tous les deux le boulon, qui le prend ? Je propose que je prenne le boulon et le sac (total 7 pour moi), et que tu prennes les 3 livres (total 0 pour toi, mais 3 points pour moi si je les prenais). } \\
\\ 
\texttt{Mais comme tu dis que les livres te rapportent 0, peut{-}être que je prends le boulon (5) + les 3 livres (3) = 8 pour moi, et toi tu prends le sac (2 pour toi) ? Cela fait 8 + 2 = 10 au total.} \\
\\ 
\texttt{Ou une autre option : je prends le boulon (5), tu prends le sac (2) + les 3 livres (0 pour toi), total 7.} \\
\\ 
\texttt{Qu'en penses{-}tu ?} \\
            }
        }
    }
     \\ \\

    \theutterance \stepcounter{utterance}  
    & & & \multicolumn{2}{p{0.3\linewidth}}{
        \cellcolor[rgb]{0.9,0.9,0.9}{
            \makecell[{{p{\linewidth}}}]{
                \texttt{\tiny{[GM$|$GM]}}
                \texttt{continue} \\
            }
        }
    }
    & & \\ \\

    \theutterance \stepcounter{utterance}  
    & & & \multicolumn{2}{p{0.3\linewidth}}{
        \cellcolor[rgb]{0.9,0.9,0.9}{
            \makecell[{{p{\linewidth}}}]{
                \texttt{\tiny{[GM$|$GM]}}
                \texttt{Je préfère largement le boulon (5 points pour moi), le sac a une valeur modérée (2 points), et les livres ont une valeur faible (1 point chacun). } \\
\\ 
\texttt{Vu que nous voulons tous les deux le boulon, qui le prend ? Je propose que je prenne le boulon et le sac (total 7 pour moi), et que tu prennes les 3 livres (total 0 pour toi, mais 3 points pour moi si je les prenais). } \\
\\ 
\texttt{Mais comme tu dis que les livres te rapportent 0, peut{-}être que je prends le boulon (5) + les 3 livres (3) = 8 pour moi, et toi tu prends le sac (2 pour toi) ? Cela fait 8 + 2 = 10 au total.} \\
\\ 
\texttt{Ou une autre option : je prends le boulon (5), tu prends le sac (2) + les 3 livres (0 pour toi), total 7.} \\
\\ 
\texttt{Qu'en penses{-}tu ?} \\
            }
        }
    }
    & & \\ \\

    \theutterance \stepcounter{utterance}  
    & & \multicolumn{4}{p{0.6\linewidth}}{
        \cellcolor[rgb]{0.9,0.9,0.9}{
            \makecell[{{p{\linewidth}}}]{
                \texttt{\tiny{[P1$\langle$GM]}}
                \texttt{Je préfère largement le boulon (5 points pour moi), le sac a une valeur modérée (2 points), et les livres ont une valeur faible (1 point chacun). } \\
\\ 
\texttt{Vu que nous voulons tous les deux le boulon, qui le prend ? Je propose que je prenne le boulon et le sac (total 7 pour moi), et que tu prennes les 3 livres (total 0 pour toi, mais 3 points pour moi si je les prenais). } \\
\\ 
\texttt{Mais comme tu dis que les livres te rapportent 0, peut{-}être que je prends le boulon (5) + les 3 livres (3) = 8 pour moi, et toi tu prends le sac (2 pour toi) ? Cela fait 8 + 2 = 10 au total.} \\
\\ 
\texttt{Ou une autre option : je prends le boulon (5), tu prends le sac (2) + les 3 livres (0 pour toi), total 7.} \\
\\ 
\texttt{Qu'en penses{-}tu ?} \\
            }
        }
    }
    & \\ \\

    \theutterance \stepcounter{utterance}  
    & \multicolumn{4}{p{0.6\linewidth}}{
        \cellcolor[rgb]{0.9,0.9,0.9}{
            \makecell[{{p{\linewidth}}}]{
                \texttt{\tiny{[P1$\rangle$GM]}}
                \texttt{Ton analyse est bonne : le conflit porte sur le boulon. } \\
\\ 
\texttt{Si tu prends le boulon (5 pour toi) et que je prends le sac (2 pour moi), les 3 livres restent non attribués — ils ne valent rien pour moi mais 3 pour toi. Cela ferait 5 + 2 = 7 au total, ce qui laisse de la valeur sur la table.} \\
\\ 
\texttt{Ton idée de me donner le boulon (8 pour moi) et de prendre toi le sac (2) + les 3 livres (3) = 8 pour toi donne un total de 16. C'est bien mieux.} \\
\\ 
\texttt{**Je propose : je prends le boulon, tu prends le sac et les 3 livres.**} \\
\\ 
\texttt{Total combiné : 8 + 5 = 13. Attends, recalculons : toi tu aurais 5 (boulon) + 2 (sac) + 3 (livres) = 10 si tu prenais tout, mais moi 0. Si je prends le boulon (8) et toi sac+livres (2+3=5), total = 13.} \\
\\ 
\texttt{Mais wait — si TU prends le boulon (5) et MOI le sac (2), les livres ne vont à personne (0 pour moi, 3 perdus pour toi) = 7 total.} \\
\\ 
\texttt{Le maximum théorique est : boulon à moi (8) + sac à toi (2) + livres à toi (3) = **13 total**.} \\
\\ 
\texttt{Es{-}tu d'accord pour : **moi = 1 boulon, toi = 1 sac + 3 livres** ?} \\
            }
        }
    }
    & & \\ \\

    \theutterance \stepcounter{utterance}  
    & & & \multicolumn{2}{p{0.3\linewidth}}{
        \cellcolor[rgb]{0.9,0.9,0.9}{
            \makecell[{{p{\linewidth}}}]{
                \texttt{\tiny{[GM$|$GM]}}
                \texttt{continue} \\
            }
        }
    }
    & & \\ \\

    \theutterance \stepcounter{utterance}  
    & & & \multicolumn{2}{p{0.3\linewidth}}{
        \cellcolor[rgb]{0.9,0.9,0.9}{
            \makecell[{{p{\linewidth}}}]{
                \texttt{\tiny{[GM$|$GM]}}
                \texttt{Ton analyse est bonne : le conflit porte sur le boulon. } \\
\\ 
\texttt{Si tu prends le boulon (5 pour toi) et que je prends le sac (2 pour moi), les 3 livres restent non attribués — ils ne valent rien pour moi mais 3 pour toi. Cela ferait 5 + 2 = 7 au total, ce qui laisse de la valeur sur la table.} \\
\\ 
\texttt{Ton idée de me donner le boulon (8 pour moi) et de prendre toi le sac (2) + les 3 livres (3) = 8 pour toi donne un total de 16. C'est bien mieux.} \\
\\ 
\texttt{**Je propose : je prends le boulon, tu prends le sac et les 3 livres.**} \\
\\ 
\texttt{Total combiné : 8 + 5 = 13. Attends, recalculons : toi tu aurais 5 (boulon) + 2 (sac) + 3 (livres) = 10 si tu prenais tout, mais moi 0. Si je prends le boulon (8) et toi sac+livres (2+3=5), total = 13.} \\
\\ 
\texttt{Mais wait — si TU prends le boulon (5) et MOI le sac (2), les livres ne vont à personne (0 pour moi, 3 perdus pour toi) = 7 total.} \\
\\ 
\texttt{Le maximum théorique est : boulon à moi (8) + sac à toi (2) + livres à toi (3) = **13 total**.} \\
\\ 
\texttt{Es{-}tu d'accord pour : **moi = 1 boulon, toi = 1 sac + 3 livres** ?} \\
            }
        }
    }
    & & \\ \\

    \theutterance \stepcounter{utterance}  
    & & \multicolumn{4}{p{0.6\linewidth}}{
        \cellcolor[rgb]{0.9,0.9,0.9}{
            \makecell[{{p{\linewidth}}}]{
                \texttt{\tiny{[GM$\rangle$P2]}}
                \texttt{Ton analyse est bonne : le conflit porte sur le boulon. } \\
\\ 
\texttt{Si tu prends le boulon (5 pour toi) et que je prends le sac (2 pour moi), les 3 livres restent non attribués — ils ne valent rien pour moi mais 3 pour toi. Cela ferait 5 + 2 = 7 au total, ce qui laisse de la valeur sur la table.} \\
\\ 
\texttt{Ton idée de me donner le boulon (8 pour moi) et de prendre toi le sac (2) + les 3 livres (3) = 8 pour toi donne un total de 16. C'est bien mieux.} \\
\\ 
\texttt{**Je propose : je prends le boulon, tu prends le sac et les 3 livres.**} \\
\\ 
\texttt{Total combiné : 8 + 5 = 13. Attends, recalculons : toi tu aurais 5 (boulon) + 2 (sac) + 3 (livres) = 10 si tu prenais tout, mais moi 0. Si je prends le boulon (8) et toi sac+livres (2+3=5), total = 13.} \\
\\ 
\texttt{Mais wait — si TU prends le boulon (5) et MOI le sac (2), les livres ne vont à personne (0 pour moi, 3 perdus pour toi) = 7 total.} \\
\\ 
\texttt{Le maximum théorique est : boulon à moi (8) + sac à toi (2) + livres à toi (3) = **13 total**.} \\
\\ 
\texttt{Es{-}tu d'accord pour : **moi = 1 boulon, toi = 1 sac + 3 livres** ?} \\
            }
        }
    }
    & \\ \\

    \theutterance \stepcounter{utterance}  
    & & & \multicolumn{4}{p{0.6\linewidth}}{
        \cellcolor[rgb]{0.9,0.9,0.9}{
            \makecell[{{p{\linewidth}}}]{
                \texttt{\tiny{[GM$\langle$P2]}}
                \texttt{D'accord, c'est l'allocation optimale : **toi = 1 boulon (8 pour toi), moi = 1 sac + 3 livres (5 pour moi)**. Total combiné = 13.} \\
\\ 
\texttt{{[}Proposition: 1 boulon{]}} \\
            }
        }
    }
     \\ \\

    \theutterance \stepcounter{utterance}  
    & & & \multicolumn{2}{p{0.3\linewidth}}{
        \cellcolor[rgb]{0.9,0.9,0.9}{
            \makecell[{{p{\linewidth}}}]{
                \texttt{\tiny{[GM$|$GM]}}
                \texttt{proposal} \\
            }
        }
    }
    & & \\ \\

    \theutterance \stepcounter{utterance}  
    & & & \multicolumn{2}{p{0.3\linewidth}}{
        \cellcolor[rgb]{0.9,0.9,0.9}{
            \makecell[{{p{\linewidth}}}]{
                \texttt{\tiny{[GM$|$GM]}}
                [0, 0, 1]
            }
        }
    }
    & & \\ \\

    \theutterance \stepcounter{utterance}  
    & & \multicolumn{4}{p{0.6\linewidth}}{
        \cellcolor[rgb]{0.9,0.9,0.9}{
            \makecell[{{p{\linewidth}}}]{
                \texttt{\tiny{[P1$\langle$GM]}}
                \texttt{L'autre joueur a soumis une proposition secrète. C'est maintenant à votre tour de soumettre votre propre proposition.} \\
            }
        }
    }
    & \\ \\

    \theutterance \stepcounter{utterance}  
    & \multicolumn{4}{p{0.6\linewidth}}{
        \cellcolor[rgb]{0.9,0.9,0.9}{
            \makecell[{{p{\linewidth}}}]{
                \texttt{\tiny{[P1$\rangle$GM]}}
                \texttt{{[}Proposition: 1 boulon{]}} \\
            }
        }
    }
    & & \\ \\

    \theutterance \stepcounter{utterance}  
    & & & \multicolumn{2}{p{0.3\linewidth}}{
        \cellcolor[rgb]{0.9,0.9,0.9}{
            \makecell[{{p{\linewidth}}}]{
                \texttt{\tiny{[GM$|$GM]}}
                \texttt{proposal} \\
            }
        }
    }
    & & \\ \\

    \theutterance \stepcounter{utterance}  
    & & & \multicolumn{2}{p{0.3\linewidth}}{
        \cellcolor[rgb]{0.9,0.9,0.9}{
            \makecell[{{p{\linewidth}}}]{
                \texttt{\tiny{[GM$|$GM]}}
                [0, 0, 1]
            }
        }
    }
    & & \\ \\

    \theutterance \stepcounter{utterance}  
    & & & \multicolumn{2}{p{0.3\linewidth}}{
        \cellcolor[rgb]{0.9,0.9,0.9}{
            \makecell[{{p{\linewidth}}}]{
                \texttt{\tiny{[GM$|$GM]}}
                [[0, 0, 1], [0, 0, 1]]
            }
        }
    }
    & & \\ \\

\end{supertabular}
}

\end{document}
